\sectionkicker{Source-backed technical notes}
\subsection*{生成から実行へ――Copilot、ReActとコード・行動・検索の接続: Technical Notes}
本欄は記事本文で圧縮した一次資料上の情報を、比較・再検証しやすい形で整理したものである。時系列、確認済みの主張、留意点、一次資料URLを掲載する。
\smallskip
\noindent{\small\textit{Technical Notesでは、一次資料から確認した公開・提供・研究上の事実を記載する。数値・能力評価や提供元・著者による比較表現は、独立再現済みの結果ではなく帰属claimとして扱う。}}\par
\medskip
\subsection*{Theme at a glance}
\addcontentsline{toc}{subsection}{Theme at a glance}
\begin{center}
\footnotesize
\begin{tabularx}{\linewidth}{@{}p{0.20\linewidth}p{0.10\linewidth}p{0.14\linewidth}X@{}}
\toprule
資料 & 位置づけ & 種別 & 時系列 \\
\midrule
GitHub Copilot general availability & 主要資料 & 製品 & 2022-06-21 \\
ReAct & 主要資料 & 論文 & 2022-10-06 \\
AlphaCode & 補足資料 & モデル & 2022-02-02, 2022-02-08 \\
Codex applications & 補足資料 & 製品 & 2022-05-24 \\
Video PreTraining (VPT) & 補足資料 & Agent & 2022-06-23 \\
Atlas & 補足資料 & 論文 & 2022-08-05 \\
\bottomrule
\end{tabularx}
\end{center}
\normalsize

\begin{technicalnote}{GitHub Copilot general availability}{主要資料}
\begingroup
% reader-facing Technical Notes break policy
\widowpenalty=10000
\clubpenalty=10000
\displaywidowpenalty=10000
\begin{tabularx}{\linewidth}{@{}>{\bfseries}p{0.22\linewidth}X@{}}
組織 & GitHub \\
種別 & 製品 \\
時系列 & 2022-06-21 \\
\end{tabularx}
\smallskip
{\bfseries 一次資料から整理したtechnical points}
\begin{itemize}[leftmargin=1.5em,itemsep=0.35em]
\item \textbf{一次情報で確認できる事実}: GitHubの選定済み一次資料では、GitHubは2022年6月21日にCopilotをindividual developers向けのgeneral availabilityへ移し、technical previewから実利用可能なcoding productへ提供状態を変更した。 製品はeditor extensionとしてcodeやcommentの文脈からreal-time coding suggestionを提示し、当時はNeovim・JetBrains IDEs・Visual Studio・Visual Studio Codeへの統合を案内していた。
\end{itemize}
\begin{minipage}{\linewidth}
% reader-facing Technical Notes generic boundary/source tail group
\begin{samepage}
{\bfseries 一次資料}
\begingroup\sloppy
\Urlmuskip=0mu plus 2mu\relax
\begin{itemize}[leftmargin=1.5em,itemsep=0.25em]
\item {\scriptsize\url{https://github.blog/news-insights/product-news/github-copilot-is-generally-available-to-all-developers/}}
\end{itemize}
\endgroup
\end{samepage}
\end{minipage}
\endgroup
\end{technicalnote}
\medskip

\begin{technicalnote}{ReAct}{主要資料}
\begingroup
% reader-facing Technical Notes break policy
\widowpenalty=10000
\clubpenalty=10000
\displaywidowpenalty=10000
\begin{tabularx}{\linewidth}{@{}>{\bfseries}p{0.22\linewidth}X@{}}
組織 & authors \\
種別 & 論文 \\
時系列 & 2022-10-06 \\
\end{tabularx}
\smallskip
{\bfseries 一次資料から整理したtechnical points}
\begin{itemize}[leftmargin=1.5em,itemsep=0.35em]
\item \textbf{一次情報で確認できる事実}: authorsの選定済み一次資料では、ReActはlanguage modelにreasoning traceとtask-specific actionをinterleaveして生成させ、推論側でaction planを追跡・更新しつつ、action側で外部knowledge baseやenvironmentから追加情報を取得する構成を提案した。 論文ではHotpotQA・FEVERでWikipedia APIとのinteractionを、ALFWorld・WebShopでinteractive decision makingを扱い、reasoningとactionを別々ではなく一つのtask-solving trajectoryとして接続した。
\end{itemize}
\begin{minipage}{\linewidth}
% reader-facing Technical Notes generic boundary/source tail group
\begin{samepage}
{\bfseries 一次資料}
\begingroup\sloppy
\Urlmuskip=0mu plus 2mu\relax
\begin{itemize}[leftmargin=1.5em,itemsep=0.25em]
\item {\scriptsize\url{http://arxiv.org/abs/2210.03629v3}}
\end{itemize}
\endgroup
\end{samepage}
\end{minipage}
\endgroup
\end{technicalnote}
\medskip

\begin{technicalnote}{AlphaCode}{補足資料}
\begingroup
% reader-facing Technical Notes break policy
\widowpenalty=10000
\clubpenalty=10000
\displaywidowpenalty=10000
\begin{tabularx}{\linewidth}{@{}>{\bfseries}p{0.22\linewidth}X@{}}
組織 & Google DeepMind \\
種別 & モデル \\
時系列 & 2022-02-02, 2022-02-08 \\
\end{tabularx}
\smallskip
{\bfseries 一次資料から整理したtechnical points}
\begin{itemize}[leftmargin=1.5em,itemsep=0.35em]
\item \textbf{一次情報で確認できる事実}: Google DeepMindの選定済み一次資料では、AlphaCodeはunseenなcompetitive programming問題に対するcode generation systemとして構成され、training/evaluation用のcleanな競技プログラミングdatasetとlarge Transformerを用いる。 解候補をlarge-scaleにsampleした後、program behaviorに基づいてfilterし、提出する少数候補へ絞り込む探索・選別pipelineを採用した。
\end{itemize}
\begin{minipage}{\linewidth}
% reader-facing Technical Notes generic boundary/source tail group
\begin{samepage}
{\bfseries 一次資料}
\begingroup\sloppy
\Urlmuskip=0mu plus 2mu\relax
\begin{itemize}[leftmargin=1.5em,itemsep=0.25em]
\item {\scriptsize\url{http://arxiv.org/abs/2203.07814v1}}
\item {\scriptsize\url{https://deepmind.google/blog/competitive-programming-with-alphacode/}}
\end{itemize}
\endgroup
\end{samepage}
\end{minipage}
\endgroup
\end{technicalnote}
\medskip

\begin{technicalnote}{Codex applications}{補足資料}
\begingroup
% reader-facing Technical Notes break policy
\widowpenalty=10000
\clubpenalty=10000
\displaywidowpenalty=10000
\begin{tabularx}{\linewidth}{@{}>{\bfseries}p{0.22\linewidth}X@{}}
組織 & OpenAI \\
種別 & 製品 \\
時系列 & 2022-05-24 \\
\end{tabularx}
\smallskip
{\bfseries 一次資料から整理したtechnical points}
\begin{itemize}[leftmargin=1.5em,itemsep=0.35em]
\item \textbf{一次情報で確認できる事実}: OpenAIの選定済み一次資料では、OpenAIは2022年5月時点でCodexをOpenAI API経由のapplication layerへ展開し、70のapplicationが複数use caseで利用している提供状態を公表した。 この記録はcode-generation研究benchmarkではなく、APIを介して外部applicationへモデル能力が接続されたproductization boundaryを示す。
\end{itemize}
\begin{minipage}{\linewidth}
% reader-facing Technical Notes generic boundary/source tail group
\begin{samepage}
{\bfseries 一次資料}
\begingroup\sloppy
\Urlmuskip=0mu plus 2mu\relax
\begin{itemize}[leftmargin=1.5em,itemsep=0.25em]
\item {\scriptsize\url{https://openai.com/index/codex-apps}}
\end{itemize}
\endgroup
\end{samepage}
\end{minipage}
\endgroup
\end{technicalnote}
\medskip

\begin{technicalnote}{Video PreTraining (VPT)}{補足資料}
\begingroup
% reader-facing Technical Notes break policy
\widowpenalty=10000
\clubpenalty=10000
\displaywidowpenalty=10000
\begin{tabularx}{\linewidth}{@{}>{\bfseries}p{0.22\linewidth}X@{}}
組織 & OpenAI \\
種別 & Agent \\
時系列 & 2022-06-23 \\
\end{tabularx}
\smallskip
{\bfseries 一次資料から整理したtechnical points}
\begin{itemize}[leftmargin=1.5em,itemsep=0.35em]
\item \textbf{一次情報で確認できる事実}: OpenAIの選定済み一次資料では、VPTは少量のlabeled Minecraft操作dataでinverse dynamics modelを学習し、それを用いて大量のunlabeled online gameplay videoへaction labelを付与してbehavioral priorを学習するsemi-supervised imitation learning構成を採る。 agentはnative mouse/keyboard interfaceでenvironmentへactionを出し、そのbehavioral priorをimitation learningやreinforcement learningでfine-tuneするため、text生成とは異なるenvironment-action layerを形成する。
\end{itemize}
\begin{minipage}{\linewidth}
% reader-facing Technical Notes generic boundary/source tail group
\begin{samepage}
{\bfseries 一次資料}
\begingroup\sloppy
\Urlmuskip=0mu plus 2mu\relax
\begin{itemize}[leftmargin=1.5em,itemsep=0.25em]
\item {\scriptsize\url{http://arxiv.org/abs/2206.11795v1}}
\item {\scriptsize\url{https://openai.com/index/vpt}}
\end{itemize}
\endgroup
\end{samepage}
\end{minipage}
\endgroup
\end{technicalnote}
\medskip

\begin{technicalnote}{Atlas}{補足資料}
\begingroup
% reader-facing Technical Notes break policy
\widowpenalty=10000
\clubpenalty=10000
\displaywidowpenalty=10000
\begin{tabularx}{\linewidth}{@{}>{\bfseries}p{0.22\linewidth}X@{}}
組織 & Meta AI / authors \\
種別 & 論文 \\
時系列 & 2022-08-05 \\
\end{tabularx}
\smallskip
{\bfseries 一次資料から整理したtechnical points}
\begin{itemize}[leftmargin=1.5em,itemsep=0.35em]
\item \textbf{一次情報で確認できる事実}: Meta AI / authorsの選定済み一次資料では、Atlasはfew-shotのknowledge-intensive taskを対象に、language modelと外部document retrievalを組み合わせてpre-trainしたretrieval-augmented language modelとして設計された。 知識をmodel parameterだけへ固定せずdocument indexから取得する構成で、論文はindex内容を更新可能な外部knowledge interfaceとして扱う。
\end{itemize}
\begin{minipage}{\linewidth}
% reader-facing Technical Notes generic boundary/source tail group
\begin{samepage}
{\bfseries 一次資料}
\begingroup\sloppy
\Urlmuskip=0mu plus 2mu\relax
\begin{itemize}[leftmargin=1.5em,itemsep=0.25em]
\item {\scriptsize\url{http://arxiv.org/abs/2208.03299v3}}
\end{itemize}
\endgroup
\end{samepage}
\end{minipage}
\endgroup
\end{technicalnote}
\medskip
